\chapter[Resultados]{Resultados: arquitectura, prototipo y comparación de estrategias}
\label{ch:resultados}

Los capítulos anteriores describieron el diseño y su marco de medida.
Este presenta su materialización y los resultados obtenidos.
Se organiza en cuatro partes: la arquitectura tecnológica que implementa la metodología (\ref{sec:res-arquitectura}), el corpus resultante (\ref{sec:res-corpus}), el prototipo \proyecto (\ref{sec:res-prototipo}) y la comparación cuantitativa de las estrategias de recuperación (\ref{sec:res-benchmark}).

\section{De la metodología a la implementación}
\label{sec:res-arquitectura}

Todas las decisiones tecnológicas comparten dos criterios heredados de los requisitos: componentes de código abierto desplegables en infraestructura propia (requisito \ref{req:soberania}) y herramientas maduras con las que el coste de operación sea asumible por una institución.
La Figura~\ref{fig:arquitectura} muestra la arquitectura completa y la Tabla~\ref{tab:stack} su correspondencia con las etapas metodológicas del Capítulo~\ref{ch:metodologia}.

\begin{figure}[htbp]
      \centering
      \resizebox{\textwidth}{!}{
            \begin{tikzpicture}[
                  font=\small,
                  node distance=7mm and 7mm,
                  proc/.style={draw=upmblue!80!black, fill=upmblue!8, thick, rounded corners=3pt, align=center, minimum height=13mm, text width=27mm, inner sep=2mm},
                  store/.style={draw=teal!60!black, fill=teal!9, thick, rounded corners=3pt, align=center, minimum height=13mm, text width=27mm, inner sep=2mm},
                  ext/.style={draw=black!60, fill=black!4, thick, rounded corners=3pt, align=center, minimum height=13mm, text width=27mm, inner sep=2mm},
                  infra/.style={draw=black!50, fill=orange!8, thick, rounded corners=3pt, align=center, minimum height=10mm, inner sep=2mm},
                  flecha/.style={-{Stealth[length=2.6mm]}, semithick, draw=black!75},
                  flechad/.style={{Stealth[length=2.6mm]}-{Stealth[length=2.6mm]}, semithick, draw=black!75}
            ]
            % Fila 1: ingesta y estructuración
            \node[ext] (portal) {\textbf{Portal Científico UPM}\\[0.5mm] \footnotesize Fuente pública};
            \node[proc, right=of portal] (descarga) {\textbf{Descarga asíncrona}\\[0.5mm] \footnotesize httpx + tenacity};
            \node[store, right=of descarga] (minio1) {\textbf{MinIO}\\[0.5mm] \footnotesize Bucket con los datos crudos};
            \node[proc, right=of minio1] (parsing) {\textbf{Parseo y modelado}\\[0.5mm] \footnotesize Scrapy + Pydantic};
            \node[store, right=of parsing] (memgraph) {\textbf{Memgraph}\\[0.5mm] \footnotesize Grafo de actividad};

            \draw[flecha] (portal) -- (descarga);
            \draw[flecha] (descarga) -- (minio1);
            \draw[flecha] (minio1) -- (parsing);
            \draw[flecha] (parsing) -- (memgraph);

            % Fila 2 (sentido inverso): enriquecimiento e indexación
            \node[proc, below=14mm of memgraph] (consol) {\textbf{Resúmenes anuales}\\[0.5mm] \footnotesize Ollama + Agno};
            \node[store, below=8mm of parsing] (minio2) {\textbf{MinIO}\\[0.5mm] \footnotesize Bucket con los perfiles enriquecidos};
            \node[proc, below=10mm of minio1] (indexa) {\textbf{Fusión e indexación}\\[0.5mm] \footnotesize Biografía vigente, dispersa y densa};
            \node[store, below=14mm of descarga] (qdrant) {\textbf{Qdrant}\\[0.5mm] \footnotesize Índices disperso y denso};

            \draw[flecha] (memgraph) -- (consol);
            \draw[flecha] (consol) -- (minio2);
            \draw[flechad] (minio2) -- (indexa);
            \draw[flecha] (indexa) -- (qdrant);

            % Fila 3: servicio en línea
            \node[ext, below=18mm of qdrant] (estudiante) {\textbf{Estudiante}\\[0.5mm] \footnotesize navegador web};
            \node[proc, below=16mm of indexa] (flask) {\textbf{\proyecto}\\[0.5mm] \footnotesize Flask + Gunicorn};
            \node[proc, below=15mm of minio2] (agentes) {\textbf{Agentes \gls{llm}}\\[0.5mm] \footnotesize Ollama (Qwen2.5)};

            \draw[flechad] (estudiante) -- (flask);
            \draw[flechad] (flask) -- (agentes);
            \draw[flechad] (flask) -- (qdrant);

            % Infraestructura transversal
            \node[infra, below=12mm of flask, text width=120mm] (infra) {\textbf{Orquestación y despliegue:} Airflow (ejecución periódica del pipeline) · Docker Compose (todos los servicios) · Traefik (proxy inverso) íntegramente \emph{on-premise}};

            \begin{scope}[on background layer]
            \node[draw=black!30, dashed, rounded corners=4pt, fill=black!1, inner sep=4mm,
                  fit=(portal)(memgraph)(qdrant)(consol),
                  label={[black!60, font=\footnotesize\bfseries]90:Preparación del corpus (fases 1--7)}] {};
            \node[draw=black!30, dashed, rounded corners=4pt, fill=black!1, inner sep=4mm,
                  fit=(estudiante)(flask)(agentes),
                  label={[black!60, font=\footnotesize\bfseries]180:Servicio}] {};
            \end{scope}
            \end{tikzpicture}
      }
      \caption{Arquitectura tecnológica del sistema.
      Las filas superiores implementan la preparación del corpus (Figura~\ref{fig:flujo}); la inferior, la atención de consultas.
      Los almacenes (MinIO, Memgraph, Qdrant) materializan los artefactos persistentes entre etapas.}
      \label{fig:arquitectura}
\end{figure}

\begin{table}[htbp]
      \centering
      \caption{Correspondencia entre etapas metodológicas y tecnologías empleadas.}
      \label{tab:stack}
      \begin{tabularx}{\textwidth}{l X l}
            \toprule
            \rowcolor{upmblue!8} \colhead{Etapa} & \colhead{Cometido} & \colhead{Tecnología} \\
            \midrule
            Adquisición & Censo de investigadores y descarga concurrente con reintentos & Python (httpx, tenacity) \\
            Capa cruda & Conservación inmutable del HTML y artefactos intermedios & MinIO \\
            Estructuración & Extracción de entidades y carga del grafo de actividad & Scrapy (selectores), Memgraph (Cypher) \\
            Validación & Esquemas tipados de perfiles y documentos & Pydantic \\
            Resúmenes anuales & Síntesis autónoma por año (resumen y dominios) con \gls{llm} local y salida estructurada & Ollama + Agno (Qwen2.5 7B) \\
            Biografía vigente & Fusión recurrente de los resúmenes anuales, ponderando el trabajo reciente & Ollama + Agno (Qwen2.5 7B) \\
            Embeddings densos & Representación semántica de la biografía vigente & nomic-embed-text-v2-moe (Ollama) \\
            Vectores dispersos & Términos ponderados por atención (temas y dominios) & BM42 (fastembed) \\
            Índice y búsqueda & Almacenamiento vectorial, similitud y fusión RRF nativa & Qdrant \\
            Aplicación web & Interfaz de búsqueda y API & Flask + Gunicorn \\
            Orquestación & Ejecución programada e incremental del pipeline & Airflow \\
            Despliegue & Contenedores y proxy inverso, todo on-premise & Docker Compose + Nginx \\
            \bottomrule
      \end{tabularx}
\end{table}

Tres decisiones merecen justificación expresa.
\duda{Creo que habría que justificar algo más cada tecnología}
Primera, el \textbf{\gls{llm} local} (Qwen2.5 7B instruccional servido por Ollama) es el punto donde el requisito \ref{req:soberania} se juega su viabilidad, pues sustituye a las \gls{api} propietarias tanto en el enriquecimiento masivo del corpus como en la interpretación de consultas en línea.
Por otro lado, Agno aporta la capa de agentes con salida estructurada validada por Pydantic.
Segunda, \textbf{Qdrant} como índice vectorial único para ambas representaciones: soporta vectores densos y dispersos en la misma colección y ofrece fusión \gls{rrf} nativa, lo que simplifica la estrategia híbrida a una sola consulta.
Además, el algoritmo BM42 fue inventado por los desarrolladores de esta tecnología, de modo que está integrado de forma nativa.
Tercera, \textbf{Airflow} como orquestador, las siete fases del pipeline son tareas idempotentes con reanudación (\emph{skip-existing}), de modo que la actualización periódica del corpus es un proceso desatendido.
Completan el cuadro \textbf{MinIO}, como capa de objetos que materializa el principio de fuente cruda inmutable y reprocesable y \textbf{Memgraph}, como base de grafo elegida por la naturaleza de vecindad del dominio (Sección~\ref{sec:met-grafo}) y porque deja preparada la línea futura de detección de colaboraciones.
El modelo de datos se mantuvo agnóstico del motor de grafo, de modo que podría migrarse a otro motor sin reescribir la lógica de carga.

\section{El corpus resultante}
\label{sec:res-corpus}

La ejecución del pipeline sobre el perímetro de las cuatro escuelas del ámbito de las Tecnologías de la Información y las Comunicaciones produce un grafo de conocimiento de 172\,621 nodos y 268\,886 relaciones.
La Tabla~\ref{tab:corpus-grafo} desglosa su composición por tipo (el esquema del grafo se detalla en el Anexo~\ref{annex:modelo-datos}).

\begin{table}[htbp]
    \centering
    \caption{Composición del grafo de conocimiento: nodos y relaciones por tipo.}
    \label{tab:corpus-grafo}
    \begin{tabularx}{\textwidth}{l r X}
        \toprule
        \rowcolor{upmblue!8} \colhead{Elemento} & \colhead{Cantidad} & \colhead{Descripción} \\
        \midrule
        \multicolumn{3}{l}{\textit{Nodos (172\,621)}} \\
        \texttt{Publication} & 109\,296 & Publicaciones (artículos, conferencias, etc.) \\
        \texttt{LastYearProject} & 44\,297  & Trabajos supervisados (TFG/TFM). \\
        \texttt{Subject} & 12\,602 & Asignaturas. \\
        \texttt{Pdi} & 5\,998 & Investigadores censados. \\
        \texttt{Position} & 428 & Unidades organizativas y cargos. \\
        \midrule
        \multicolumn{3}{l}{\textit{Aristas (268\,886)}} \\
        \texttt{AUTHORSHIP} & 157\,556 & Autorías ($\mathrm{Pdi}\rightarrow\mathrm{Publication}$). \\
        \texttt{TUTORING} & 48\,916 & Supervisiones ($\mathrm{Pdi}\rightarrow\mathrm{LastYearProject}$). \\
        \texttt{TEACHING} & 32\,567 & Docencia ($\mathrm{Pdi}\rightarrow\mathrm{Subject}$). \\
        \texttt{COAUTHORED} & 19\,125  & Coautorías entre investigadores ($\mathrm{Pdi}\rightarrow\mathrm{Pdi}$). \\
        \texttt{WORKS} & 10\,722 & Pertenencia a unidades ($\mathrm{Pdi}\rightarrow\mathrm{Position}$). \\
        \bottomrule
    \end{tabularx}
\end{table}

Las 428 unidades organizativas (\texttt{Position}) heredan las categorías del portal: 226 grupos de investigación, 88 áreas de conocimiento, 62 departamentos universitarios, 22 institutos o centros de investigación, 20 escuelas y un puñado de unidades de otro tipo.

La afiliación de los investigadores a las escuelas merece una advertencia, pues condiciona la lectura de estas cifras y la propia delimitación del corpus.
El portal científico no enlaza de forma sistemática a cada investigador con su escuela: de los 5\,998 investigadores del grafo, solo 1\,431 tienen una relación directa con un nodo de tipo escuela, mientras que 4\,567 carecen por completo de ella. En particular, 2\,216 están vinculados a un departamento pero no a la escuela correspondiente.
Es una limitación de la fuente, no del modelo: ese enlace sencillamente no figura en muchas fichas.

Esta carencia obliga a acotar el perímetro de forma indirecta.
Como no puede filtrarse de manera fiable por escuela (la mayoría de los investigadores cuelgan del departamento y no de la escuela), el censo se delimita a través de los departamentos, lo que es determinista pero imperfecto debido a la peculiar organización interna de la \gls{upm}: algunos departamentos imparten en varias escuelas aunque \emph{oficialmente} pertenezcan a una sola (por ejemplo, el Departamento de Lingüística Aplicada a la Ciencia y la Tecnología, o el de Ingeniería de Organización, Administración de Empresas y Estadística, adscrito a la Escuela de Industriales pero con profesorado que ha impartido docencia en otras escuelas).
En consecuencia, el criterio puede incorporar a algún investigador que enseña en otra escuela y omitir a otros que sí lo hacen en las del perímetro.
Se asume como mal menor, destinado a diluirse cuando el corpus se extienda al conjunto de la universidad.

Aplicando ese criterio determinista, el perímetro efectivo de las cuatro escuelas (sus cuatro nodos de escuela y catorce departamentos adscritos) reúne 996 investigadores distintos; los 5\,998 nodos \texttt{Pdi} del grafo incluyen, además, a coautores externos arrastrados por las relaciones de actividad.
La Tabla~\ref{tab:corpus-escuelas} desglosa el perímetro por escuela, junto con la cobertura de temas declarados en el portal.
Veinticuatro investigadores están adscritos a unidades de más de una escuela, por lo que la suma por escuela supera ligeramente el total de investigadores distintos.

\begin{table}[htbp]
    \centering
    \caption{Investigadores del perímetro por escuela y cobertura de temas declarados.}
    \label{tab:corpus-escuelas}
    \begin{tabularx}{\textwidth}{X r r}
        \toprule
        \rowcolor{upmblue!8} \colhead{Escuela} & \colhead{Investigadores} & \colhead{Con temas declarados} \\
        \midrule
        \gls{etsii}  & 315 & 188 (59,7\,\%) \\
        \gls{etsit}  & 310 & 287 (92,6\,\%) \\
        \gls{etsisi} & 237 & 125 (52,7\,\%) \\
        \gls{etsist} & 158 & 132 (83,5\,\%) \\
        \midrule
        Total (investigadores distintos) & 996 & 709 (71,2\,\%) \\
        \bottomrule
    \end{tabularx}
\end{table}

De ese perímetro, los \textbf{709} \glspl{pdi} con contenido aprovechable (los mismos que declaran temas en la Tabla~\ref{tab:corpus-escuelas}) constituyen la población efectivamente indexada.
La colección dispersa \texttt{researcher\_\allowbreak keywords} (fase 5) está completa con esos 709 puntos, uno por \gls{pdi}.
La colección \texttt{researcher\_\allowbreak summaries} (fase 7) mantiene un punto por \gls{pdi} y, al concluir, contiene un total de 754.
\pendiente{Esta colección al fina tiene 754 PDIs, tengo que mirar por qué y explicitarlo en algún lado}
Conviene contrastar este perímetro con la dimensión total del censo: el grafo contiene 5\,998 nodos \texttt{Pdi}, de modo que la indexación actual cubre una fracción del profesorado de la universidad y deja margen para la extensión al conjunto de la \gls{upm}.

\subsection{Tiempos del pipeline}
\label{sec:res-tiempos}

La Tabla~\ref{tab:tiempos-fases} recoge la duración aproximada de cada fase en una ejecución completa sobre el perímetro, realizada en el servidor de despliegue.

\begin{table}[htbp]
    \centering
    \caption{Duración aproximada de cada fase del pipeline en una ejecución completa sobre el perímetro (servidor de despliegue).}
    \label{tab:tiempos-fases}
    \begin{tabularx}{\textwidth}{l X r}
        \toprule
        \rowcolor{upmblue!8} \colhead{Fase} & \colhead{Cometido} & \colhead{Duración aprox.} \\
        \midrule
        1. Censo de identificadores & Descubrimiento de investigadores en el portal & $<$ 10 s \\
        2. Descarga & Volcado del perfil y la actividad pública & $\sim$ 57 min \\
        3. Estructuración & Parseo y carga del grafo en Memgraph & $>$ 4 h 30 min \\
        4. Enriquecimiento & Consolidación del perfil & $>$ 11 min \\
        5. Indexación de keywords & Colección dispersa \texttt{researcher\_keywords} (solo para los \glspl{pdi} seleccionados) & $\sim$ 25 min \\
        6. Resúmenes anuales & Síntesis por año con \gls{llm} local (solo para los \glspl{pdi} seleccionados) & 29 h ($\approx$ 1 día y 5 h) \\
        7. Fusión e indexación & Biografía vigente y \texttt{researcher\_summaries} (solo para los \glspl{pdi} seleccionados) & \ $\sim$ 7h 30 min \\
        \bottomrule
    \end{tabularx}
\end{table}

El reparto de costes es muy desigual y deja una lectura nítida.
Las fases de ingesta y estructuración (1-4), más la indexación léxica (5), se completan en unas pocas horas, dominadas por el parseo y la carga del grafo (fase 3).
El coste decisivo es la \emph{generación de los resúmenes con el modelo de lenguaje local}.
Tan solo la fase 6 se prolonga alrededor de un día y cinco horas sobre los 995 \glspl{pdi} y la fase 7 (fusión recurrente e indexación) añade todavía más tiempo.
Y ello pese a emplear un modelo deliberadamente ligero (Qwen2.5-7B-Instruct) y sin modo de razonamiento explícito, elegido precisamente para maximizar el rendimiento de generación.
Como las fases 5 y 6 se ejecutan en paralelo (Sección~\ref{sec:res-arquitectura}), la indexación de keywords queda absorbida dentro de la ventana de los resúmenes.

Que la preparación del corpus se mida en días, y no en minutos, no es un detalle de implementación sino una restricción de diseño: es lo que convierte el procesamiento por lotes periódico, desatendido e incremental (\emph{skip-existing}) en una necesidad y no en una comodidad (\ref{req:adquisicion}).
Esto justifica que la orquestación con Airflow reanude solo lo nuevo o lo inválido en lugar de repetir el proceso completo.
La corrida completa se ejecutó en el servidor de despliegue, no en una estación de trabajo local.
Las implicaciones de esta dependencia del hardware para la escalabilidad del proyecto (y el coste de procesar en el futuro toda la \gls{upm}) se analizan en la Sección~\ref{sec:disc-coste-llm}.

\pendiente{Completar al concluir la fase 7: su duración (Tabla~\ref{tab:tiempos-fases}) y la cobertura de biografía y de dominios por escuela.}

\section{El prototipo \proyecto}
\label{sec:res-prototipo}

El prototipo materializa la capa de presentación descrita en la Sección~\ref{sec:met-presentacion}: una aplicación web con un único campo de texto en el que el estudiante describe su idea en cualquier idioma, un selector de estrategia (con la híbrida por defecto) y una lista de resultados ordenada.

Cada tarjeta de resultado muestra el nombre del investigador, \textcolor{red}{su puntuación normalizada respecto al primer resultado}, los temas de su perfil que coinciden con la consulta y, en las estrategias semánticas, su biografía investigadora.
Cuando interviene el \gls{llm}, la interfaz muestra también la interpretación realizada (los términos extraídos o la reformulación enriquecida), de modo que el estudiante pueda auditar y, en su caso, reformular su consulta.
La latencia percibida varía según la estrategia: las que incorporan el \gls{llm} añaden el coste de una generación local a la búsqueda vectorial.
Las puramente vectoriales son inmediatas (la semántica directa resuelve en torno a dos décimas de segundo), mientras que las que invocan el modelo en consulta se sitúan entre los 7 y los 19 segundos de media (Tabla~\ref{tab:benchmark-tiempos}), un coste dominado por la generación local.

\pendiente{Incluir una o dos capturas de pantalla del prototipo (búsqueda y resultados, idealmente con la misma consulta en dos modos distintos) y un párrafo guiando su lectura.}

El prototipo está desplegado y es accesible públicamente en \url{https://wiig.dia.fi.upm.es/upm-match/}, donde cualquier usuario puede introducir una idea de trabajo y obtener recomendaciones en tiempo real.

\section{Comparación de estrategias de recuperación}
\label{sec:res-benchmark}

La Tabla~\ref{tab:benchmark} resume los resultados del banco de pruebas descrito en el Capítulo~\ref{ch:evaluacion} para las cinco estrategias.
Las consultas son las propuestas reformuladas en primera persona (estándar de plata, Sección~\ref{sec:eval-groundtruth}).
De los 500 trabajos del conjunto de referencia, 458 quedan tras restringirlas al universo común de investigadores presentes en todos los índices (Sección~\ref{sec:eval-protocolo}) y sobre ese conjunto se calculan todas las métricas.

\begin{table}[htbp]
      \centering
      \caption[Resultados del banco de pruebas sobre la verdad de referencia de supervisión histórica]{Resultados del banco de pruebas sobre la verdad de referencia de supervisión histórica. Las columnas son las métricas que reporta el arnés de evaluación, por haber un único relevante por consulta, $\mathrm{MAP}$ coincide con la $\mathrm{MRR}$ (Sección~\ref{sec:eval-metricas}).}
      \label{tab:benchmark}
      \footnotesize
      \setlength{\tabcolsep}{4pt}
      \begin{tabularx}{\textwidth}{X c c c c c c c c}
            \toprule
            \rowcolor{upmblue!8} \colhead{Estrategia} & \colhead{MRR} & \colhead{MAP} & \colhead{Hit@1} & \colhead{Hit@3} & \colhead{Hit@5} & \colhead{Hit@10} & \colhead{nDCG@5} & \colhead{nDCG@10} \\
            \midrule
            Léxica con \gls{llm} & 0,091 & 0,091 & 0,041 & 0,094 & 0,146 & 0,218 & 0,093 & 0,116 \\
            Léxica directa (control) & 0,063 & 0,063 & 0,033 & 0,066 & 0,098 & 0,135 & 0,065 & 0,077 \\
            Semántica directa & 0,137 & 0,137 & 0,072 & 0,166 & 0,197 & 0,266 & 0,139 & 0,161 \\
            Semántica enriquecida & 0,146 & 0,146 & 0,085 & 0,168 & 0,205 & 0,284 & 0,148 & 0,173 \\
            Híbrida & \textbf{0,153} & \textbf{0,153} & \textbf{0,090} & \textbf{0,175} & \textbf{0,218} & \textbf{0,303} & \textbf{0,156} & \textbf{0,184} \\
            \bottomrule
      \end{tabularx}
\end{table}

El tiempo medio por consulta (Tabla~\ref{tab:benchmark-tiempos}) cuantifica el coste de incorporar el \gls{llm} al bucle de consulta, separando el coste de generación (extracción o enriquecimiento) del de la búsqueda vectorial propiamente dicha.

\begin{table}[htbp]
      \centering
      \caption[Tiempo medio por consulta de cada estrategia]{Tiempo por consulta de cada estrategia, en segundos, sobre la misma máquina e índices ($n=458$). La columna \textit{LLM} es la fracción del tiempo total dedicada a la generación (extracción de términos o enriquecimiento); \textit{búsqueda} es el resto (embebido y consulta a Qdrant).}
      \label{tab:benchmark-tiempos}
      \footnotesize
      \begin{tabularx}{\textwidth}{X c c c c c}
            \toprule
            \rowcolor{upmblue!8} \colhead{Estrategia} & \colhead{Media} & \colhead{Mediana} & \colhead{p90} & \colhead{LLM} & \colhead{Búsqueda} \\
            \midrule
            Semántica directa & 0,23  & 0,21  & 0,22  & 0,00  & 0,23 \\
            Léxica directa (control) & 2,57  & 2,50  & 3,30  & 0,00  & 2,57 \\
            Léxica con \gls{llm} & 7,59  & 7,50  & 9,86  & 5,99  & 1,60 \\
            Semántica enriquecida & 11,51 & 11,42 & 12,68 & 7,61  & 3,90 \\
            Híbrida & 18,69 & 18,50 & 20,60 & 13,15 & 5,54 \\
            \bottomrule
      \end{tabularx}
\end{table}

De estas cifras se extraen cuatro lecturas, todas ellas a interpretar como \emph{comparación relativa} entre estrategias y no como medidas absolutas de utilidad (Sección~\ref{sec:eval-validez}).

Primera, \textbf{la señal semántica domina con holgura a la léxica}.
Las tres estrategias densas superan a las dos léxicas en todas las métricas y cortes: la peor de las semánticas (la directa, \gls{mrr} 0,137) aventaja con claridad a la mejor de las léxicas (la guiada por \gls{llm}, 0,091) y esta a su vez al control de coincidencia léxica pura (0,063).
La afinidad temática entre la idea del estudiante y la trayectoria del investigador se captura, por tanto, mucho mejor en el espacio denso de los \emph{embeddings} que mediante la coincidencia de términos.

Segunda, \textbf{el \gls{llm} en consulta aporta en ambos frentes}.
En el léxico, extraer términos con el modelo (0,091) mejora en cerca de la mitad sobre el control directo (0,063): la reinterpretación rescata coincidencias que el texto crudo no producía.
Y en el semántico, el enriquecimiento (0,146) supera ahora a la semántica directa (0,137) en \gls{mrr} y en todos los cortes, a diferencia de lo observado con el texto sin reformular.
Al partir de una propuesta breve y en lenguaje natural, la traducción y expansión controlada del modelo añade señal recuperable que el embebido directo no alcanzaba.

Tercera, \textbf{la hibridación es la mejor estrategia y justifica su diseño}.
La híbrida encabeza \emph{todas} las métricas de la Tabla~\ref{tab:benchmark} (\gls{mrr} 0,153, \mbox{Hit@1} 0,090, \mbox{Hit@10} 0,303), por delante del enriquecimiento y de la semántica directa.
Fusionar la señal dispersa explicable con la densa no solo amplía la cobertura en los cortes bajos, sino que mejora el acierto en cabeza, lo que respalda la decisión de adoptarla como estrategia por defecto del prototipo (Sección~\ref{sec:res-prototipo}).

Cuarta, \textbf{la calidad se paga en latencia}.
La estrategia ganadora es también la más cara: la híbrida resuelve en 18,7\,s de media (Tabla~\ref{tab:benchmark-tiempos}), dominados por la generación local, frente a los 0,23\,s de la semántica directa, que no invoca el \gls{llm} en consulta.
Existe, pues, un compromiso real entre calidad y coste (Sección~\ref{sec:disc-coste-llm}): la semántica directa es una alternativa atractiva cuando la latencia sea crítica, a cambio de una pérdida moderada de acierto (\gls{mrr} 0,137 frente a 0,153).

Conviene subrayar que esta jerarquía solo emerge con las propuestas reformuladas.
En una corrida exploratoria previa, tomando el texto de referencia directamente como consulta, el orden se invertía: la semántica directa lideraba y ni el enriquecimiento ni la hibridación aportaban nada.
El motivo es el solapamiento léxico entre el texto original del trabajo y el perfil de su propio tutor, que premia artificialmente a las estrategias más literales (la amenaza de optimismo por solapamiento de la Sección~\ref{sec:eval-validez}).
La reformulación en primera persona reduce ese solapamiento y revela el valor real del \gls{llm} y de la fusión de señales, lo que constituye, además, una validación empírica de la conveniencia de construir las consultas mediante el estándar de plata.

El desglose por tipo de trabajo, escuela y curso (Anexo~\ref{annex:resultados-extendidos}) matiza estas tendencias: el liderazgo de las estrategias enriquecida e híbrida es transversal, con la \gls{etsit} como la escuela mejor recuperada y la \gls{etsii} la más difícil.
Entre \glspl{tfg} y \glspl{tfm} las diferencias son pequeñas.

\section{Casos ilustrativos}
\label{sec:res-casos}

\pendiente{Incluir dos o tres casos cualitativos end-to-end: la consulta de ejemplo (una en español coloquial, otra técnica en inglés), los tres primeros investigadores devueltos por la estrategia híbrida con sus evidencias, y un breve comentario sobre la pertinencia de cada recomendación.
Elegir casos que ilustren tanto un acierto claro como un fallo instructivo, en línea con la discusión del Capítulo~\ref{ch:discusion}.}
